% ===================================================================
%  పట్టిక సంఖ్య 10 — సముద్రం. ఓడరేవు.
%
%  Traduction, faite en 2026, en telougou d'aujourd'hui, sur l'ido de
%  texto/io. Regles de la colonne : voir 10-pattika-01.tex. Une seule
%  scene, vingt-cinq alineas, cent renvois : les cles ne portent ni
%  c1 ni c2.
%
%  LE TITRE DU TABLEAU TIENT DANS LE BLOC « apar », sous le numero,
%  comme aux tableaux 1 a 6 et comme le font deja les colonnes hindi
%  et marathe.
%
%  LE FAC-SIMILE IDO GRASSE « kande » au bloc 06-1 : c'est un
%  mot-outil, et les colonnes traduites ne le suivent pas la-dessus
%  (voir le degrassage des mots-outils, deja fait partout). Il grasse
%  en revanche une phrase entiere a la fin du bloc 05-1 — celle-la se
%  garde, comme en hindi.
%
%  CE TABLEAU-CI, LE TELOUGOU L'ATTENDAIT. Les huit precedents
%  butaient sur une chose absente tous les deux ou trois tableaux —
%  le marron, le fleau, le hetre. Ici rien ne manque : l'Andhra a
%  neuf cents kilometres de cote, un port de guerre a Visakhapatnam,
%  et la langue a les mots, y compris les recents :
%
%    జలాంతర్గామి (10) le sous-marin — mot telougou, pas un emprunt
%    కవచ నౌక (9)   le cuirasse     చుక్కాని (30) le gouvernail
%    లంగరు (77)    l'ancre         తెడ్డు (47)   l'aviron
%    దీపస్తంభం (93) le phare       ద్వీపకల్పం (91) la presqu'ile
%    జలసంధి (98)   le detroit      అఖాతాలు (95)  les baies
%    పోటుపాటు      la maree        వల (56)       le filet
%    తెరచాప (13)   la voile        గిడ్డంగి (43) l'entrepot
%
%  Restent les emprunts d'un vocabulaire technique qui s'est fait en
%  anglais partout, telougou compris, et que le telougou d'aujourd'hui
%  dit en anglais : క్రూయిజర్ (6), ఫ్రిగేట్ (11), టార్పెడో (8),
%  బాయిలర్లు (25), డాకులు (26), ప్రొపెల్లర్ (28), క్రేన్లు (62),
%  లారీలు (64), ట్రాము (74), కాసినో (89). Ce ne sont pas des trous
%  comme le marron chaud : ce sont les mots que la langue emploie.
%  Ecrire un calque savant a leur place aurait ete la seule fausse
%  note du tableau.
%
%  SEIZE INVERSIONS SUR LES TRENTE PAIRES DE parigi.py, et le tri des
%  quatorze autres est celui du tableau 9, aux memes trois raisons :
%  enumeration (« kesti (60) od en pakegi (61) », « komandanto (83)
%  o dil admiralo (84) »), renvoi de gauche deja modificateur
%  (« tramvoyo (74) di qua la halteyo (75) » — le telougou dit
%  « l'arret du tram » dans cet ordre-la sans qu'on touche a rien), et
%  frontiere de phrase (« infantrio (80) kun sua sabro. Dop li esis
%  navano (81) »).
%
%  ET UNE INVERSION QUE LA LISTE N'A PAS VUE, PAR PORTEE. « ponteto
%  (20), qua komunikigas la parad-esplanado kun l'arsenalo (21) » :
%  le relatif est bien la, mais 55 caracteres separent les deux
%  renvois — mesures, pas estimes — et PORTEO vaut 45. C'est le
%  premier des deux echappements que parigi.py declare dans son
%  entete : celui de longue distance, deja rencontre au tableau 5.
%  On n'allonge pas la portee pour autant, parce que la table de
%  calibrage de l'entete dit qu'a 80 le rappel ne monte plus et que
%  le nombre de lignes a lire, lui, passe de 55 a 88. La liste fait
%  gagner du temps, pas de la certitude : c'est renvoji.py qui
%  tranche, et il a tranche — zero divergence.
% ===================================================================

%%K t10-apar-1 apar
\VUsaut{4.0mm}
\VUtitre{15.0pt}{60}{పట్టిక సంఖ్య 10}
\VUsaut{3.0mm}
\VUpk{11.4pt}{40}{సముద్రం. --- ఓడరేవు.}
\VUsaut{3.0mm}

%%K t10-01-1 p
1. --- వేసవిలో కొందరు కొండపై ప్రశాంతతా చల్లదనమూ వెతుక్కుంటే, ఇంకొందరు తీరానికి వెళ్ళడమే
మేలనుకుంటారు. మేము గత సంవత్సరం అదే చేశాము — ఎందుకంటే మాకు \VUgras{మహాసముద్రం}
\textsuperscript{(1)} అంటే చాలా ఇష్టం.

%%K t10-02-1 p
2. --- నా మటుకు, \VUgras{క్షితిజం} \textsuperscript{(2)} దాకా పరుచుకున్న ఆ అపారమైన నీటి
విస్తీర్ణాన్ని చూస్తూ ఉంటే నాకు గొప్ప ఆనందం; అమెరికాకు వెళ్ళే ప్రయాణికులు ఎక్కే భారీ
\VUgras{ఆవిరి ఓడలు} \textsuperscript{(3)} సుదీర్ఘ \VUgras{సముద్రయానానికి} బయలుదేరడం
చూసినప్పుడూ అంతే.

%%K t10-03-1 p
3. --- అందుకే, నా అభిరుచి తెలిసిన మా నాన్న, సెలవులు గడిపేందుకు ఎప్పుడూ మన పెద్ద
\VUgras{నౌకాదళ రేవుల్లో} ఒకదాని దగ్గరున్న ప్రశాంతమైన ఇసుక తీరాన్నే ఎంచుకుంటారు.

%%K t10-03-2 p
మేము చివరిసారి అక్కడ ఉన్నప్పుడు, \VUgras{నౌకాదళ శ్రేణి} భారీ \VUgras{అడ్డుకట్ట}
\textsuperscript{(4)} వెనుక లంగరు వేసింది — ఆ కట్టే \VUgras{నౌకాదళ రేవును}
\textsuperscript{(5)} మూసి కాపాడుతుంది —, నేను కోరుకున్నట్టే రకరకాల \VUgras{యుద్ధ నౌకలను}
తీరిగ్గా చూడగలిగాను : నీళ్ళ కిందికి మునిగి మాయమయ్యే చిన్న \VUgras{జలాంతర్గామిని}
\textsuperscript{(10)}, ఇంకా భారీ \VUgras{కవచ నౌకను} \textsuperscript{(9)} — ఇప్పటిదాకా
దానికి భయం అంటే దాదాపు \VUgras{టార్పెడో పడవ} \textsuperscript{(8)} దాడులంటేనే.

%%K t10-tit-2 sub
\VUsaut{5.0mm}
\VUpk{10.2pt}{40}{చిన్న నౌకలు.}
\VUsaut{3.4mm}

%%K t10-04-1 p
4. --- మందపాటి లోహ కవచంలో బలహీనమైన చోటును ఎంతో \VUgras{నేర్పుగా} కనిపెట్టడం ఆ పడవకు
అప్పటికే తెలుసు — అక్కడ ప్రమాదకరమైన \VUgras{టార్పెడోను} అమర్చేందుకు; దాని పేలుడు నౌక
నాశనానికి కారణమవుతుంది. అయినా జలాంతర్గామితో పోలిస్తే దానికి భయపడాల్సింది తక్కువ, ఎందుకంటే
\VUgras{టార్పెడో వ్యతిరేక నౌకలు} దాన్ని సులభంగా తరుముతాయి.

%%K t10-05-1 p
5. --- వేగవంతమైన కవచపు \VUgras{క్రూయిజర్ నౌకలను} \textsuperscript{(6)} కూడా నేను చూడగలిగాను
— నిజమైన కవచ నౌకంత దాదాపు గాయపడనివి —, కొన్ని \VUgras{రవాణా నౌకలను}
\textsuperscript{(12)}, మూడు నాలుగు \VUgras{ఫిరంగి పడవలను} \textsuperscript{(7)}, చివరికి ఒక
\VUgras{ఫ్రిగేట్‌ను} \textsuperscript{(11)}.
\VUgras{ఆ నౌకా శ్రేణి లంగరు ఎత్తేందుకు కదులుతున్నప్పటి దృశ్యం అన్నిటికంటే ఆసక్తికరం.}

%%K t10-06-1 p
6. --- ఆ పెద్ద ఉక్కు కుప్పలకూ, అందమైన \VUgras{మూడు స్తంభాల ఓడలకూ}, ఇప్పటికీ వ్యాపార
నౌకాశ్రేణిలో వాడుకలో ఉన్న ముద్దైన \VUgras{తెరచాప ఓడలకూ} \textsuperscript{(13)} నిర్మాణంలో
ఎంత తేడా! నేను \VUgras{రేవు దిమ్మెపై} \textsuperscript{(14)} నడుస్తున్నప్పుడు, తెరచాపలన్నీ
విప్పుకుని అవి రేవులోకి రావడం చూసేవాణ్ణి. నిజమే, \VUgras{గాలి} ఎదురైనప్పుడు వాటికున్న మేలు
చాలావరకూ పోయేది. కుంటలోకి తిరిగి రావాలంటే తెరచాపలన్నీ దించుకోవాల్సి వచ్చేది,
\VUgras{లాగుడు నౌక} \textsuperscript{(16)} ఒకటి కావాల్సి వచ్చేది — వాటిని తీసుకెళ్ళి,
మామూలు \VUgras{సరుకు పడవల} \textsuperscript{(17)} లాగే రేవు గట్టు దాకా తేవడానికి.

%%K t10-tit-3 sub
\VUsaut{5.0mm}
\VUpk{10.2pt}{40}{వ్యాపార రేవు.}
\VUsaut{3.4mm}

%%K t10-07-1 p
7. --- నేను చెబుతున్న ఈ రేవులో ఆసక్తికరమైనది ఏమిటంటే — భారీ పనులతో కృత్రిమంగా తీర్చిదిద్దిన
నౌకాదళ రేవు మాత్రమే కాదు, పెద్ద నదీ \VUgras{ముఖద్వారంలో} ఉన్న అద్భుతమైన
\VUgras{వ్యాపార రేవు} కూడా ఇందులో ఉంది.

%%K t10-08-1 p
8. --- ఆ రెండు కుంటలను వేరుచేసే భూమి నాలికకు కోటబందీ ఉంది; సముద్రంలోకి చొచ్చుకుపోయే
\VUgras{ఫిరంగి మెట్లూ} \VUgras{బురుజులూ} దాన్ని కాపాడతాయి. \VUgras{కోట గోడలపై}
\textsuperscript{(15)} నడవాలంటే ప్రత్యేక అనుమతి కావాలి. మా మామయ్య ద్వారా అది నాకు సులభంగానే
దొరికింది — ఆయన \VUgras{ఓడల నిర్మాణశాలల} \textsuperscript{(18)} ఇంజనీరు —, పెద్ద పాకల కింద
కడుతున్న నౌకలను చూసేందుకు వెళ్ళాను.

%%K t10-09-1 p
9. --- ఒక కవచ నౌకను నీళ్ళలోకి దించడం కూడా నేను చూశాను; ఆ పెద్ద ముద్ద తనకు తానే
వదిలేయబడి, ఊయల లాంటి చట్రంపై జారడం మొదలుపెట్టి, నదీ నీటిపై తేలడం మొదలుపెట్టిన ఆ క్షణంలో
జనమంతా పొందిన ఉద్వేగం నాకు గుర్తుంది.

%%K t10-10-1 p
10. --- నిర్మాణశాలలకు వెళ్ళాలంటే \VUgras{కవాతు మైదానం} \textsuperscript{(19)} దాటాలి — అక్కడ
దండులోని సైనికులు రోజూ వ్యాయామం చేసేందుకు వస్తారు —, ఆ తర్వాత ఇనుప \VUgras{చిన్న వంతెన}
\textsuperscript{(20)} మీదుగా వెళ్ళాలి; అది కవాతు మైదానాన్ని \VUgras{ఆయుధాగారంతో}
\textsuperscript{(21)} కలుపుతుంది.

%%K t10-tit-4 sub
\VUsaut{5.0mm}
\VUpk{10.2pt}{40}{ఆయుధాగారం, డాకులు.}
\VUsaut{3.4mm}

%%K t10-11-1 p
11. --- మా మామయ్యతో కలిసి నేను ఆ పెద్ద కట్టడాన్ని చూశాను — అది \VUgras{ఫిరంగులు}
\textsuperscript{(22)}, \VUgras{ఫిరంగి గుళ్ళు} \textsuperscript{(23)}, \VUgras{పేలుడు గుళ్ళతో}
నిండి ఉంది. \VUgras{లోహపు కర్మాగారాన్ని} \textsuperscript{(24)} కూడా చూశాను — అక్కడే ఆవిరి
ఓడల \VUgras{బాయిలర్లు} \textsuperscript{(25)} తయారవుతాయి.

%%K t10-12-1 p
12. --- చాలా దగ్గరలోనే \VUgras{డాకులు} \textsuperscript{(26)} — మరమ్మతు డాకులు — ఉన్నాయి,
ఇంకా చాలా విశేషమైన \VUgras{తేలియాడే డాకులు} \textsuperscript{(27)} ఉన్నాయి; వాటిలో మరమ్మతుకు
వచ్చిన ఒక నౌకపై పడవ \VUgras{ప్రొపెల్లరును} \textsuperscript{(28)}, \VUgras{అడుగు దూలాన్ని}
\textsuperscript{(29)}, \VUgras{చుక్కానిని} \textsuperscript{(30)} చాలా దగ్గరగా చూడగలిగాను.

%%K t10-13-1 p
13. --- నౌకాదళ రేవంత గానే వ్యాపార రేవూ చూడదగ్గది; రేవు గట్లపై నేను చాలా గంటలు నోరు
తెరుచుకుని నిలబడ్డాను — అక్కడ నదిని ఎదురీదే \VUgras{తెడ్డు చక్రాల ఓడలు}
\textsuperscript{(31)} కట్టేసి ఉంటాయి —, \VUgras{స్తంభాలెత్తే క్రేన్లు}
\textsuperscript{(32)} పనిచేయడం చూస్తూ; అవి భారీ సరుకులను ఈకల్లా ఎత్తుతాయి.

%%K t10-tit-5 sub
\VUsaut{5.0mm}
\VUpk{10.2pt}{40}{నౌకా మార్గదర్శి.}
\VUsaut{3.4mm}

%%K t10-14-1 p
14. --- రేవు దాటి బయటికి వెళ్ళే \VUgras{అట్లాంటిక్ ఓడలను} \textsuperscript{(34)} చూసి
మురిసిపోయాను — వాటి \VUgras{జెండాలు} \textsuperscript{(35)} గాలిలో ఎగురుతున్నాయి, నేర్పరి
\VUgras{నౌకా మార్గదర్శి} \textsuperscript{(36)} వాటిని నడిపిస్తున్నాడు; \VUgras{తేలే గుర్తులతో},
\VUgras{దారి కర్రలతో} \textsuperscript{(40)} గుర్తుపెట్టిన సన్నని \VUgras{కాలువ} వెంట
నడపడం అతనికి అద్భుతంగా తెలుసు — ఒకే ఒక చదరపు తెరచాపతో కష్టంగా దారి తీసుకునే
\VUgras{చదునైన సరుకు పడవలను} \textsuperscript{(41)} తప్పించుకుంటూ. \VUgras{ముందు భాగంలో}
\textsuperscript{(37)} — ముక్కులో — రెండో తరగతి \VUgras{గదులను} \textsuperscript{(39)}, ఇంకా
\VUgras{వెనుక భాగంలో} \textsuperscript{(38)} — తోకలో — మొదటి తరగతి ప్రయాణికుల విడిది గదులనూ
నేను గుర్తుపట్టాను.

%%K t10-15-1 p
15. --- \VUgras{సరుకు పడవలో} \textsuperscript{(42)} సరుకు ఎక్కించడం కూడా అప్పుడప్పుడూ
చూశాను — అందులో \VUgras{గిడ్డంగి} \textsuperscript{(43)} సరుకూ, \VUgras{ఓడరేవు స్టేషను}
\textsuperscript{(44)} సరుకూ ఇప్పుడే కుప్పబోశారు; \VUgras{రేవు గట్టు రైలుదారి}
\textsuperscript{(45)} మీద వచ్చే బోలెడన్ని రైళ్ళు వాటిని తెచ్చిపెట్టాయి.

%%K t10-tit-6 sub
\VUsaut{5.0mm}
\VUpk{10.2pt}{40}{తెడ్డు విహారం.}
\VUsaut{3.4mm}

%%K t10-16-1 p
16. --- \VUgras{నీటి డాకులో} \textsuperscript{(53)} నేను తరచూ పడవ నడిపేవాణ్ణి — అందులోకే
\VUgras{పడవలు} \textsuperscript{(46)} ఆశ్రయం కోసం వస్తాయి —, \VUgras{తెడ్డు}
\textsuperscript{(47)} వేయడం నాకు ఆనందంగా ఉండేది. \VUgras{డెక్కు} \textsuperscript{(48)} మీదికి
ఎక్కాను — \VUgras{తెరచాప పడవది} \textsuperscript{(49)} అది —, \VUgras{తాళ్ళ}
\textsuperscript{(52)} సాయంతో \VUgras{స్తంభం} \textsuperscript{(51)} ఎక్కి
\VUgras{అడ్డకొయ్యల} \textsuperscript{(50)} దాకా వెళ్ళాను; ఆ తర్వాత \VUgras{సన్నని పడవలో}
\textsuperscript{(54)} నా విహారం కొనసాగించాను — బరువైన \VUgras{చేపల పడవల}
\textsuperscript{(55)} మధ్య, వాటిపైనే \VUgras{వలలు} \textsuperscript{(56)} ఆరబెట్టి ఉంటాయి.

%%K t10-17-1 p
17. --- \VUgras{రేవు గట్టుపై} \textsuperscript{(58)} నా కళ్ళముందు రకరకాల \VUgras{సరుకులు}
\textsuperscript{(59)} వరుసకట్టాయి — \VUgras{పెట్టెల్లో} \textsuperscript{(60)}, లేదా
\VUgras{మూటల్లో} \textsuperscript{(61)} ఉన్నవి. అవి పెద్ద పడవల \VUgras{సరుకు మోత}
\textsuperscript{(63)} అయ్యాయి; బలమైన \VUgras{క్రేన్లు} \textsuperscript{(62)} వాటిని దించి
\VUgras{లారీలపై} \textsuperscript{(64)} పెట్టాయి, ఈలోగా \VUgras{సరుకు రవాణాదారులు} వాటిని
ప్రకటించి \VUgras{సుంకపు కార్యాలయంలో} \textsuperscript{(65)} పన్ను కట్టారు.

%%K t10-18-1 p
18. --- చివరికి \VUgras{తిరిగే వంతెన} \textsuperscript{(66)} దగ్గరికో \VUgras{నీటి తలుపు}
\textsuperscript{(69)} దగ్గరికో వచ్చినప్పుడు — \VUgras{బురద తీసే యంత్రం}
\textsuperscript{(68)} ముందునుంచి వెళ్తూ, అది \VUgras{కాలువను} \textsuperscript{(67)}
శుభ్రం చేస్తుంది —, \VUgras{సంకేత స్తంభం} \textsuperscript{(70)} దగ్గర రేవు దిమ్మెపై నేను
దిగాను.

%%K t10-19-1 p
19. --- ఆ దిమ్మెకు అవతలి వైపునే \VUgras{నౌకా పడవలు} \textsuperscript{(71)} వచ్చి ఆగుతాయి —
అవి నౌకాదళ శ్రేణి అధికారులను ఒడ్డుకు తెస్తాయి. నావికులు తాళానికి తగ్గట్టు తమ తెడ్లు ఎత్తడం
చూడటం ఒక అద్భుత దృశ్యం; ఈలోగా \VUgras{అల} \textsuperscript{(72)} మోసుకొచ్చిన పడవ దిగుడు మెట్ల
దగ్గర సులభంగా ఆగుతుంది. ఈ చోటినుంచే \VUgras{ఇసుక తీరంపై} \textsuperscript{(73)}
\VUgras{పోటుపాటునూ}, \VUgras{పోటు}, \VUgras{పాటుల} కదలికనూ బాగా గమనించవచ్చు.

%%K t10-tit-7 sub
\VUsaut{5.0mm}
\VUpk{10.2pt}{40}{అధికారుల క్లబ్బు.}
\VUsaut{3.4mm}

%%K t10-20-1 p
20. --- ఇంటికి తిరిగి రావడానికి నేను \VUgras{విద్యుత్ ట్రాము} \textsuperscript{(74)}
ఎక్కాను; దాని \VUgras{ఆగే చోటు} \textsuperscript{(75)} అధికారుల క్లబ్బు ముందు పాతిన
\VUgras{స్తంభం} దగ్గరే ఉంది.

%%K t10-20-2 p
క్లబ్బు \VUgras{బాల్కనీలోంచి} \textsuperscript{(76)} — దాన్ని \VUgras{లంగర్లతో}
\textsuperscript{(77)}, ఫిరంగులతో, ఫిరంగి గుళ్ళతో అలంకరించారు — నౌకాదళ అధికారుల అద్భుతమైన
గుంపును నేను తరచూ చూశాను : \VUgras{లెఫ్టినెంట్లు} \textsuperscript{(78)}, తమ
\VUgras{ఖడ్గాలతో}; \VUgras{ఫిరంగిదళ కెప్టెన్లు} \textsuperscript{(79)},
\VUgras{పదాతిదళ కెప్టెన్లు} \textsuperscript{(80)}, తమ \VUgras{వంపు కత్తులతో}. వాళ్ళ వెనుక
\VUgras{నావికుడు} \textsuperscript{(81)} ఒకడు ఉండేవాడు — దూరంగా జరుగుతున్నది వాళ్ళు
చూడాలనుకున్నప్పుడు \VUgras{దూరదర్శినిని} అందించేవాడు.

%%K t10-21-1 p
21. --- కుర్ర \VUgras{ఉప లెఫ్టినెంట్లను} \textsuperscript{(82)} కూడా నేను చూశాను — వాళ్ళు
తమ \VUgras{కమాండరు} \textsuperscript{(83)} నుంచో \VUgras{అడ్మిరల్} \textsuperscript{(84)}
నుంచో ఆదేశాలు అందుకుంటున్నారు; ఈలోగా \VUgras{ముఖ్య నావికుడు} \textsuperscript{(85)} వాటిని
నౌకకు తీసుకెళ్ళేందుకు ఎదురుచూస్తున్నాడు.

%%K t10-22-1 p
22. --- ఆ బాల్కనీలోంచి కనిపించే దృశ్యం అద్భుతం : \VUgras{గుడారాలు}
\textsuperscript{(86)}, \VUgras{స్నానపు గదులతో} \textsuperscript{(87)} సహా ఇసుక తీరమంతా కంటికి
అందుతుంది; స్నానాల వేళ \VUgras{స్నానికులు} \textsuperscript{(88)} \VUgras{కాసినో}
\textsuperscript{(89)} ముందు సందడిగా గంతులేయడం కనిపిస్తుంది.

%%K t10-23-1 p
23. --- ఇసుక తీరం చివరన తీరరేఖ చిన్న \VUgras{ద్వీపకల్పంగా} \textsuperscript{(91)}
\VUgras{రాళ్ళతో} మారుతుంది — దానిమీదికే \VUgras{పెద్ద అల} \textsuperscript{(92)} వచ్చి
విరుగుతుంది, దానిమీదే \VUgras{దీపస్తంభం} \textsuperscript{(93)} కట్టారు; రాత్రిపూట అది
నావికులకు దారి చూపుతుంది. ఆ ద్వీపకల్పానికి నేలతో ఉన్న సంబంధం చాలా సన్నని
\VUgras{సంధి భూమి} \textsuperscript{(90)} ద్వారానే.

%%K t10-tit-8 sub
\VUsaut{5.0mm}
\VUpk{10.2pt}{40}{తీరాల సాధారణ దృశ్యం.}
\VUsaut{3.4mm}

%%K t10-24-1 p
24. --- \VUgras{కొండగోడ} \textsuperscript{(94)} పరుచుకుని ఉంది, సముద్రం దాన్ని చీల్చేసింది;
ఆ సముద్రమే దానిలో ఇష్టం వచ్చినట్టు వరుసగా \VUgras{అఖాతాలనూ} \textsuperscript{(95)},
\VUgras{అగ్రాలనూ} (కేపులు) గీసి, దాన్ని అత్యంత రమణీయంగా చేసింది.

%%K t10-24-2 p
\VUgras{పీఠభూమిపై} \textsuperscript{(96)} — అది \VUgras{జలసంధిని} \textsuperscript{(98)}
పైనుంచి కాచి ఉంటుంది — చాలా ముఖ్యమైన \VUgras{కోట} \textsuperscript{(97)} ఒకటి, కొన్ని
«విల్లా»లూ ఉన్నాయి.

%%K t10-25-1 p
25. --- ఆ జలసంధి ఖండం నుంచి చాలా ప్రమాదకరమైన \VUgras{ద్వీపాన్ని} \textsuperscript{(99)}
వేరుచేస్తుంది; దాని చుట్టూ \VUgras{నీటి అడుగు బండలూ} \textsuperscript{(100)},
\VUgras{అలలు విరిగే చోట్లూ} ఉన్నాయి, అక్కడ పడవలూ, పెద్ద నౌకలూ కూడా అప్పుడప్పుడూ
ఒడ్డెక్కి ఇరుక్కుంటాయి. సాయం వెంటనే అందినా, \VUgras{ప్రాణరక్షక పడవలు} వేగంగా వచ్చినా, ఆ
\VUgras{నౌకా విధ్వంసాలు} భయంకరమైనవి; విషువత్తు నాటి \VUgras{తుఫానుల్లో} చాలా నౌకలు
మనుషులతో, సరుకుతో సహా మునిగిపోయాయి.

%%K t10-25-2 p
అంత దగ్గరి ప్రమాదం ఉన్నా, \VUgras{ఆ రేవుకు} నావికులు బాగా వస్తుంటారు.

\VUsaut{6.0mm}
\VUfilet{22mm}
